\documentclass[11pt,letterpaper]{article}
\usepackage[T1]{fontenc}
\usepackage[utf8]{inputenc}
\usepackage{lmodern}
\usepackage[margin=0.95in]{geometry}
\usepackage{amsmath,amssymb,amsthm}
\usepackage{enumitem}
\usepackage{microtype}
\usepackage{xcolor}
\usepackage{booktabs,array}
\usepackage{pict2e}
\usepackage{fancyhdr}
\usepackage{hyperref}
\definecolor{mathjockgreen}{HTML}{355439}
\hypersetup{colorlinks=true,linkcolor=mathjockgreen,urlcolor=mathjockgreen,
  pdftitle={Lay Linear Algebra and Its Applications: worked solutions},
  pdfauthor={Mathjock},pdfsubject={Detailed worked solutions, fourth edition}}
\setlength{\parindent}{0pt}
\setlength{\parskip}{0.55em}
\setlength{\headheight}{14pt}
\setlist[enumerate]{leftmargin=*,itemsep=0.4em,topsep=0.4em}
\setlist[itemize]{leftmargin=*,itemsep=0.25em}
\allowdisplaybreaks
\emergencystretch=2em
% Leave a visible gap after two-digit exercise numbers in the contents.
\makeatletter
% Chapter.section and supplementary labels need more room than article's
% default 1.5em section-number box, especially for labels such as 3.S.
\renewcommand*{\l@section}[2]{\begingroup\bfseries
  \@dottedtocline{1}{0em}{3em}{#1}{#2}\endgroup}
\renewcommand*{\l@subsection}{\@dottedtocline{2}{1.5em}{3.5em}}
\makeatother
\pagestyle{fancy}
\fancyhf{}
\fancyhead[L]{\small Mathjock \textbar\ Lay, fourth edition}
\fancyhead[R]{\small Step-by-step solutions}
\fancyfoot[C]{\thepage}
\newcommand{\N}{\mathbb{N}}
\newcommand{\Z}{\mathbb{Z}}
\newcommand{\Q}{\mathbb{Q}}
\newcommand{\R}{\mathbb{R}}
\newcommand{\abs}[1]{\left|#1\right|}
\makeatletter
\newcommand{\exercise}[2]{\subsection*{Exercise #1: #2}\addcontentsline{toc}{subsection}{Exercise #1: #2}\def\@currentlabel{#1}\label{ex:#1}}
\makeatother
\newcommand{\goal}{\par\textbf{What we must show.} }
\newcommand{\idea}{\par\textbf{Plan.} }
\newcommand{\step}[1]{\par\medskip\textbf{#1}\par}
\newcommand{\answer}[1]{\par\textbf{Conclusion.} #1\par}
\newcommand{\C}{\mathbb{C}}
\DeclareMathOperator{\Span}{Span}
\DeclareMathOperator{\Col}{Col}
\DeclareMathOperator{\Nul}{Nul}
\DeclareMathOperator{\rank}{rank}
\DeclareMathOperator{\dist}{dist}
\newcommand{\practice}[2]{\subsection*{Practice problem #1: #2}\addcontentsline{toc}{subsection}{Practice problem #1: #2}}
\newcommand{\bookinfo}{%
\textbf{Text and numbering.} David C. Lay, \emph{Linear Algebra and Its
Applications}, fourth edition (Pearson, 2012), as supplied in the textbook
PDF linked on Mathjock. Labels give chapter, section, and exercise number;
for example, 1.2.3 is Exercise 3 in Section 1.2. A P label denotes a practice
problem, and S denotes a chapter supplementary exercise.
These are independently written solutions, not a publisher's solutions manual.
Short descriptions identify the problems; consult the textbook for the full
original statements.\par
\textbf{Source scope.} The supplied textbook contains printed Chapters~1--8.
Its separately published online Chapters~9--10 are absent from that PDF and
are not included in this collection.\par
\textbf{How to read the solutions.} Start with the section's definitions.
For a calculation, follow the row operations and substitute the result into
the original equations. For a proof, identify the assumptions, the desired
conclusion, and why each implication follows. Free variables are parameters,
not values that should be silently set to zero. Unless otherwise specified,
scalars are real; complex arithmetic is used where the problem requires it.
}
